\chapter{Limits and troubleshooting}\label{ch:limits}
\chapterintro{What this version of AgenticTO does not do, what the paper's evidence does not establish, and what to do
when something goes wrong.}

\section{Scope of the software}\label{sec:scope}
\begin{itemize}
  \item \textbf{One physics model.} Isotropic, small-displacement, linear-elastic compliance minimization under one
  volume constraint, with one load case: all loads act together. No stress limits, buckling, fatigue, contact,
  nonlinear materials, multiple load cases or manufacturing constraints.
  \item \textbf{Box domains on a regular voxel grid.} No import of existing geometry (STL or STEP) as a design domain;
  complex envelopes are described with keep-out regions.
  \item \textbf{Element-unit filters.} The density filter radius is given in elements, so the smallest member shrinks in
  physical size as the grid is refined. \code{min\_feature} only feeds a warning (\rref{R009}); it does not impose a
  physical length scale.
  \item \textbf{GPU memory bounds the release.} The default ladder stops below \LadderMax\ elements. In the paper the
  certified releases were 12.3--73.7 million elements in the session study and 25.6--73.7 million in the release-policy
  study; a guarded $1.13\times10^{8}$-element attempt did not fit beside a normal desktop's GPU memory use \claim{C3.2}.
  \item \textbf{Interaction set.} The workspace has no brush editing, design variants, comparison views, parameter sweeps
  or automatic sketch on every edit. It has been tested in Chrome on Windows with one user per session.
  \item \textbf{Agents.} The built-in assistant uses Codex; other agents connect through MCP. How much clarification
  helps depends on whether the model asks (Section~\ref{sec:x2}).
  \item \textbf{Platform.} Tested on Windows 11 with an NVIDIA RTX 4090. The spill guard reads Windows GPU-memory
  counters.
\end{itemize}

\section{What the evidence does not establish}\label{sec:validity}
The paper states these limits; they apply to every use of the numbers in this book.
\begin{itemize}
  \item The 50 reference readings and all 41 contested responses were adjudicated by two authors, not by independent
  engineers; they agreed on 49 of 50 readings and 33 of 41 contested responses. Accuracy counts and paired $p$
  values stay provisional until an adjudication outside the author team exists.
  \item The suite is small and was developed alongside the system. Only two providers, one model each, were tested.
  \item A certificate is a set of numerical acceptance checks on one computation. It is not a structural safety
  assessment, and the exported surface needs further analysis before anything is built.
  \item Warm releases reach different local optima from cold runs. Neither warm policy guarantees the cold optimum or new
  members; a duct-bridge problem gave no usable comparison at all.
  \item Timings are single-workstation measurements on a shared desktop GPU. Backend latency excludes browser rendering,
  and the interaction model excludes waiting, errors and learning. No participant study has been run.
  \item Replay shows that recorded decisions reproduce recorded results. It does not show that a fresh agent run would
  decide the same way.
\end{itemize}

\section{Troubleshooting}\label{sec:trouble}
\index{troubleshooting}

\begin{reftable}{@{}>{\raggedright\arraybackslash}p{0.32\linewidth}X@{}}
Symptom & Cause and fix \\ \midrule
The \gui{GPU} pill stays \emph{starting} or turns \emph{dead}. & The worker could not load CuPy or the solver. Check
\code{AGENTTO\_VOXFEM\_SRC} and the CUDA 12 installation, then read the backend console. \\
The ladder is missing or greyed out. & The specification has an error, or is still incomplete. Fix what \gui{Checks}
lists in red. \\
A release stops with a spill message. & Another program holds GPU memory and the solve would spill into shared memory.
Close GPU-heavy programs (browsers with hardware acceleration, games, other CUDA jobs), or release a coarser rung with
\code{--ladder-max}. \\
A release is \gui{not certified}. & Read the \fail{} on the card. \emph{converged} failing means the iteration cap was
reached: try \gui{Optimize release} or check the specification. \emph{spill monitored} failing means the result is not
on the release rung. A volume or keep-in/out failure points to a problem worth reporting. \\
The first sketch of a new specification takes several seconds. & Expected: the first solve builds the multigrid
hierarchy and runs the cold schedule (2.5--6.7~s in the latency study). Later sketches after an edit are warm. \\
A rung takes far longer than its neighbours. & The grid may be padded for multigrid (rule \rref{R012} may say so) or
the problem converges slowly; preview and refine can reach their iteration cap. \gui{Stop} cancels the solve. \\
The release grid is finer than the specification's $h$. & By design: the release rung is the finest that fits on the
GPU. \tool{ladder} explains this in a note. The certificate's \code{regularization} record gives both element edges
and the filter length at each, so the shorter physical filter of the release is on record. \\
Rule \rref{R013} on a large rung. & The memory estimate exceeds the GPU. Use a coarser rung, or cap the ladder. \\
The assistant refuses to start: \code{PolicyError}. & An API-key variable is set, or the CLI is not on a subscription
login. Unset the variable or log in again (Section~\ref{sec:agents}). \\
The assistant answers but changes nothing. & In Propose mode changes come back as cards: accept them. In Act mode, open
the steps drawer and look for failed calls. \\
The assistant asks about details the brief already states. & Known behaviour (Section~\ref{sec:x2}). Answer, or choose
the reading that matches the brief. \\
An external agent's calls are all rejected. & For Codex set \code{default\_tools\_approval\_mode = "approve"}; for
Claude Code allow the \code{mcp\_\_agentto} tools. \\
An agent cancels and resubmits a long solve. & Its tool call timed out. Tell it to call \tool{wait_job}; raise the
client's tool timeout. \\
\tool{solve} from a script fails with ``spec is not complete/valid''. & Set the volume fraction and the resolution with
\tool{set_params} first (Section~\ref{sec:python}). \\
The workspace shows \gui{live session} instead of b06. & The backend session already holds a design (yours or an
agent's). Pick a task from the problem picker, or restart the backend for an empty session. \\
The front end cannot reach the backend. & The backend listens on \code{127.0.0.1:8765} and the front end proxies
\code{/api} and \code{/ws} to it; start the backend first and open \url{http://127.0.0.1:5173}. \\
\end{reftable}

\section{Reporting a problem}
Include the session folder (at least \code{session.jsonl} and the report of the affected design), the output of the
backend console, the AgenticTO version and the GPU model. A log lets the authors replay exactly what happened. The
repository's issue tracker will be at \url{\repourl} once the code is public.
